\section{Conclusion}
\label{sec:conclusion}

Agreement between two model judges is useful only to the extent that their
errors differ. By measuring internal probe readouts, we find that shared error
remains common even when the models come from different families. The
decomposition of false agreement shows why. Judge accuracy and error overlap
must be considered jointly, and improving one does not guarantee improvement in
the other. False agreement can be detected within a domain, but its signature
does not transfer uniformly to unseen data distributions. Cross-model oversight
should therefore measure shared error directly rather than infer independence
from model family or accuracy alone.
